% !TEX spellcheck = es-ES

% Plan de Trabajo de Tesis - Maestría en Ciencia de Datos
\documentclass[11pt,a4paper]{article}

% Packages
\usepackage[utf8]{inputenc}
\usepackage[spanish]{babel}
\usepackage[margin=1in]{geometry}
\usepackage[style=apa,backend=biber]{biblatex}
\addbibresource{referencias.bib}
\usepackage{booktabs}
\usepackage{longtable}
\usepackage{array}
\usepackage{hyperref}
\usepackage{xcolor}
\usepackage{listings}
\usepackage{graphicx}
\usepackage{fancyhdr}
\usepackage{amssymb}
\usepackage{amsmath}
\usepackage{titlesec}
\usepackage{float}
\usepackage{caption}
\usepackage{subcaption}

% Hyperref setup
\hypersetup{
    colorlinks=true,
    linkcolor=blue,
    filecolor=magenta,
    urlcolor=cyan,
    pdftitle={Plan de Trabajo de Tesis - Maestría en Ciencia de Datos},
    pdfauthor={Franco Astegiano},
}

% Page style
\pagestyle{fancy}
\fancyhf{}
\rhead{Plan de Tesis}
\lhead{Maestría en Ciencia de Datos}
\cfoot{\thepage}

% Code listings setup
\lstset{
    language=Python,
    basicstyle=\ttfamily\small,
    keywordstyle=\color{blue},
    commentstyle=\color{gray},
    stringstyle=\color{red},
    showstringspaces=false,
    breaklines=true,
    frame=single,
    numbers=left,
    numberstyle=\tiny\color{gray}
}

% Title setup
\title{\textbf{Plan de Trabajo de Tesis}\\
\Large{Maestría en Ciencia de Datos}\\
\vspace{0.5cm}
\large{[Título de la Tesis]}}
\author{Estudiante: Franco Astegiano\\
Director: [Nombre del Director]\\
Codirector: [Nombre del Codirector (si aplica)]}
\date{\today}

\begin{document}

\maketitle

\newpage
\tableofcontents
\newpage

% =============================================================================
\section{Información General}

\subsection{Datos del Estudiante}
\begin{itemize}
    \item \textbf{Nombre completo:} Franco Astegiano
    \item \textbf{Carrera:} Maestría en Ciencia de Datos
    \item \textbf{Universidad:} [Nombre de la Universidad]
    \item \textbf{Correo electrónico:} [correo@ejemplo.com]
    \item \textbf{Fecha de inicio del plan:} [Fecha]
\end{itemize}

\subsection{Información del Director}
\begin{itemize}
    \item \textbf{Nombre:} [Nombre completo]
    \item \textbf{Institución:} [Institución]
    \item \textbf{Correo electrónico:} [correo@ejemplo.com]
    \item \textbf{Área de especialización:} [Área]
\end{itemize}

\subsection{Información del Codirector (si aplica)}
\begin{itemize}
    \item \textbf{Nombre:} [Nombre completo]
    \item \textbf{Institución:} [Institución]
    \item \textbf{Correo electrónico:} [correo@ejemplo.com]
    \item \textbf{Área de especialización:} [Área]
\end{itemize}

% =============================================================================
\section{Resumen}



% =============================================================================
\section{Introducción}

\subsection{Contexto y Motivación}
[Descripción del contexto general del problema y la motivación para realizar esta investigación.]

\subsection{Planteamiento del Problema}
[Descripción clara y precisa del problema que se abordará en la tesis.]

\subsection{Justificación}
[Explicación de la importancia y relevancia del problema, tanto desde el punto de vista académico como práctico.]

\subsection{Alcance}
[Definición clara de lo que incluirá y no incluirá el proyecto de tesis.]

% =============================================================================
\section{Marco Teórico}

\subsection{Conceptos Fundamentales}
[Definición de los conceptos clave necesarios para entender el proyecto.]

\subsection{Estado del Arte}
[Revisión de la literatura existente y trabajos relacionados.]

\subsection{Herramientas y Tecnologías}
[Descripción de las herramientas, tecnologías y frameworks que se utilizarán.]

% =============================================================================
\section{Objetivos}

\subsection{Objetivo General}
[Objetivo principal de la tesis, debe ser claro, específico y alcanzable.]

\subsection{Objetivos Específicos}
\begin{enumerate}
    \item [Objetivo específico 1]
    \item [Objetivo específico 2]
    \item [Objetivo específico 3]
    \item [Objetivo específico 4]
    \item [Objetivo específico 5]
\end{enumerate}

% =============================================================================
\section{Metodología}

\subsection{Enfoque Metodológico}
[Descripción del enfoque general que se utilizará para abordar el problema.]

\subsection{Fases del Proyecto}

\subsubsection{Fase 1: [Nombre de la Fase]}
\begin{itemize}
    \item \textbf{Descripción:} [Descripción detallada]
    \item \textbf{Actividades:}
    \begin{enumerate}
        \item [Actividad 1]
        \item [Actividad 2]
        \item [Actividad 3]
    \end{enumerate}
    \item \textbf{Entregables:} [Lista de entregables]
    \item \textbf{Duración estimada:} [X semanas/meses]
\end{itemize}

\subsubsection{Fase 2: [Nombre de la Fase]}
\begin{itemize}
    \item \textbf{Descripción:} [Descripción detallada]
    \item \textbf{Actividades:}
    \begin{enumerate}
        \item [Actividad 1]
        \item [Actividad 2]
        \item [Actividad 3]
    \end{enumerate}
    \item \textbf{Entregables:} [Lista de entregables]
    \item \textbf{Duración estimada:} [X semanas/meses]
\end{itemize}

\subsubsection{Fase 3: [Nombre de la Fase]}
\begin{itemize}
    \item \textbf{Descripción:} [Descripción detallada]
    \item \textbf{Actividades:}
    \begin{enumerate}
        \item [Actividad 1]
        \item [Actividad 2]
        \item [Actividad 3]
    \end{enumerate}
    \item \textbf{Entregables:} [Lista de entregables]
    \item \textbf{Duración estimada:} [X semanas/meses]
\end{itemize}

\subsubsection{Fase 4: [Nombre de la Fase]}
\begin{itemize}
    \item \textbf{Descripción:} [Descripción detallada]
    \item \textbf{Actividades:}
    \begin{enumerate}
        \item [Actividad 1]
        \item [Actividad 2]
        \item [Actividad 3]
    \end{enumerate}
    \item \textbf{Entregables:} [Lista de entregables]
    \item \textbf{Duración estimada:} [X semanas/meses]
\end{itemize}

\subsection{Datos y Fuentes de Información}
[Descripción de los datos que se utilizarán, sus fuentes y métodos de recolección.]

\subsection{Herramientas y Tecnologías}
[Descripción detallada de las herramientas de software, librerías, frameworks y plataformas que se utilizarán.]

\subsection{Métricas de Evaluación}
[Definición de las métricas que se utilizarán para evaluar los resultados del proyecto.]

% =============================================================================
\section{Cronograma}

\subsection{Planificación Temporal}

\begin{table}[H]
\centering
\begin{tabular}{|l|c|c|c|c|c|c|}
\hline
\textbf{Actividad} & \textbf{Mes 1} & \textbf{Mes 2} & \textbf{Mes 3} & \textbf{Mes 4} & \textbf{Mes 5} & \textbf{Mes 6} \\
\hline
Revisión bibliográfica & X & X & & & & \\
\hline
Recolección de datos & & X & X & & & \\
\hline
Preprocesamiento & & & X & X & & \\
\hline
Desarrollo del modelo & & & & X & X & \\
\hline
Evaluación y validación & & & & & X & X \\
\hline
Redacción de tesis & & & X & X & X & X \\
\hline
Preparación defensa & & & & & & X \\
\hline
\end{tabular}
\caption{Cronograma general del proyecto de tesis}
\end{table}

\subsection{Hitos Principales}
\begin{enumerate}
    \item \textbf{[Nombre del Hito 1]:} [Descripción] - [Fecha]
    \item \textbf{[Nombre del Hito 2]:} [Descripción] - [Fecha]
    \item \textbf{[Nombre del Hito 3]:} [Descripción] - [Fecha]
    \item \textbf{[Nombre del Hito 4]:} [Descripción] - [Fecha]
    \item \textbf{[Nombre del Hito 5]:} [Descripción] - [Fecha]
\end{enumerate}

% =============================================================================
\section{Recursos Necesarios}

\subsection{Recursos Humanos}
\begin{itemize}
    \item \textbf{Estudiante:} Dedicación de [X] horas semanales
    \item \textbf{Director:} [X] horas mensuales de supervisión
    \item \textbf{Codirector:} [X] horas mensuales de supervisión (si aplica)
    \item \textbf{Otros colaboradores:} [Descripción si aplica]
\end{itemize}

\subsection{Recursos Tecnológicos}
\begin{itemize}
    \item \textbf{Hardware:} [Descripción de computadoras, servidores, GPU, etc.]
    \item \textbf{Software:} [Lista de software y licencias necesarias]
    \item \textbf{Infraestructura:} [Cloud computing, almacenamiento, etc.]
\end{itemize}

\subsection{Recursos Financieros}
\begin{table}[H]
\centering
\begin{tabular}{|l|r|}
\hline
\textbf{Concepto} & \textbf{Costo Estimado} \\
\hline
Licencias de software & \$[XX, XXX] \\
\hline
Servicios cloud & \$[XX, XXX] \\
\hline
Bibliografía & \$[XX, XXX] \\
\hline
Otros gastos & \$[XX, XXX] \\
\hline
\textbf{Total} & \textbf{\$[XX, XXX]} \\
\hline
\end{tabular}
\caption{Presupuesto estimado del proyecto}
\end{table}

\subsection{Acceso a Datos}
[Descripción de cómo se obtendrá acceso a los datos necesarios, incluyendo permisos y consideraciones éticas.]

% =============================================================================
\section{Resultados Esperados}

\subsection{Contribuciones Científicas}
[Descripción de las contribuciones que se espera hacer al campo de estudio.]

\subsection{Productos Entregables}
\begin{enumerate}
    \item \textbf{Documento de Tesis:} [Descripción]
    \item \textbf{Código Fuente:} [Descripción del repositorio y documentación]
    \item \textbf{Datasets Procesados:} [Descripción]
    \item \textbf{Modelos Entrenados:} [Descripción]
    \item \textbf{Publicaciones:} [Artículos o conferencias planeadas]
    \item \textbf{Otros:} [Aplicaciones, dashboards, APIs, etc.]
\end{enumerate}

\subsection{Impacto Esperado}
[Descripción del impacto académico, práctico o social que se espera del proyecto.]

% =============================================================================
\section{Riesgos y Mitigación}

\subsection{Identificación de Riesgos}

\begin{table}[H]
\centering
\begin{tabular}{|p{4cm}|p{2cm}|p{2cm}|p{5cm}|}
\hline
\textbf{Riesgo} & \textbf{Probabilidad} & \textbf{Impacto} & \textbf{Estrategia de Mitigación} \\
\hline
[Riesgo 1] & Alta/Media/Baja & Alto/Medio/Bajo & [Estrategia] \\
\hline
[Riesgo 2] & Alta/Media/Baja & Alto/Medio/Bajo & [Estrategia] \\
\hline
[Riesgo 3] & Alta/Media/Baja & Alto/Medio/Bajo & [Estrategia] \\
\hline
[Riesgo 4] & Alta/Media/Baja & Alto/Medio/Bajo & [Estrategia] \\
\hline
\end{tabular}
\caption{Matriz de riesgos del proyecto}
\end{table}

\subsection{Plan de Contingencia}
[Descripción de acciones alternativas en caso de que los riesgos se materialicen.]

% =============================================================================
\section{Aspectos Éticos y Legales}

\subsection{Consideraciones Éticas}
[Discusión de las consideraciones éticas relacionadas con el uso de datos, privacidad, sesgo algorítmico, etc.]

\subsection{Protección de Datos}
[Descripción de cómo se protegerán los datos sensibles y se cumplirá con regulaciones como GDPR, etc.]

\subsection{Propiedad Intelectual}
[Clarificación sobre la propiedad del código, datos y resultados del proyecto.]

% =============================================================================
\section{Plan de Comunicación}

\subsection{Reuniones de Seguimiento}
[Frecuencia y formato de las reuniones con el director y co-director.]

\subsection{Informes de Avance}
[Calendario de informes de progreso que se entregarán.]

\subsection{Presentaciones}
[Presentaciones intermedias o en conferencias que se planean realizar.]

\subsection{Publicaciones}
[Plan para publicar resultados en revistas o conferencias.]

% =============================================================================
\section{Referencias}

% Las referencias se gestionan en el archivo referencias.bib
% Para citar en el texto usa: \cite{clave}, \textcite{clave}, \parencite{clave}
% Ejemplos:
% - Según \textcite{autor2023}, el modelo...
% - Varios estudios \parencite{autor2023,autor2024} muestran...
% - Como se menciona en la literatura \cite{autor2023}...

\printbibliography[heading=none]

% =============================================================================
\section{Anexos}

\subsection{Anexo A: Glosario de Términos}
[Definición de términos técnicos utilizados en el documento.]

\subsection{Anexo B: Detalles Técnicos Adicionales}
[Información técnica complementaria que no se incluyó en el cuerpo principal.]

\subsection{Anexo C: Cartas de Apoyo}
[Cartas de instituciones o empresas que proporcionarán datos o recursos.]

% =============================================================================
\newpage
\section*{Firmas de Aprobación}

\vspace{2cm}

\begin{center}
\begin{tabular}{p{6cm} p{6cm}}
\hrulefill & \hrulefill \\
Firma del Estudiante & Fecha \\
Franco Astegiano & \\
& \\
& \\
\hrulefill & \hrulefill \\
Firma del Director & Fecha \\
[Nombre del Director] & \\
& \\
& \\
\hrulefill & \hrulefill \\
Firma del Co-director & Fecha \\
[Nombre del Co-director] & \\
\end{tabular}
\end{center}

\end{document}
